\subsection{Linear Representations}

DEFINITION 1. --- Let M be a K-module. A linear representation of G in M is a group homomorphism from G to the linear group GL(M) (II, §1, No.~2, p.~195).

Let \(\pi\) be a linear representation of G in a K-module M. We also say that the pair \((M,\pi)\) is a linear representation of G. When K is nonzero and M is a free K-module, the dimension of M is called the degree (or dimension) of the representation \(\pi\) .

Recall that the canonical mapping \(g \mapsto e_{g}\) from G to the algebra K{[}G{]} of the group G is a basis of the K-module K{[}G{]}. By abuse of notation, we also denote the element \(e_{g}\) of K{[}G{]} by g. Given a linear representation \(\pi : G \to \mathbf{GL}(M)\) of G, there exists a unique homomorphism of K-algebras from K{[}G{]} to \(\operatorname{End}_{\mathrm{K}}(\mathrm{M})\) that extends \(\pi\) (III, §2, No.~6, p.~447, Example); we also denote this extension by \(\pi\) . It is a representation of the algebra K{[}G{]} in M (VIII, p.~373, Definition 1), so that M is a K{[}G{]}-module. Conversely, every linear representation \(\rho\) of the algebra K{[}G{]} in M defines a linear representation of G in M given by \(g \mapsto \rho(g)\) .

We will freely apply the terminology concerning the K{[}G{]}-module structure to linear representations of the group G. For example, we speak of a subrepresentation, a simple representation, a direct sum of representations, etc. Let \((\mathrm{M},\pi)\) and \((\mathrm{M}^{\prime},\pi^{\prime})\) be linear representations of G; we denote by \(\operatorname{Hom}_{\mathrm{G}}(\pi,\pi^{\prime})\) the K-module \(\operatorname{Hom}_{\mathrm{K[G]}}(\mathrm{M},\mathrm{M}^{\prime})\) of homomorphisms of K{[}G{]}-modules from M to M'. \(^{(1)}\)

A mapping \(f\) from \(G\) to \(K\) is called a central function if we have \(f(gg') = f(g'g)\) for every pair \((g, g')\) of elements of \(G\) ; equivalently, we have \(f(ghg^{-1}) = f(h)\) for every pair of elements \(g, h\) of \(G\) . The central functions are therefore the functions on \(G\) whose restriction to each conjugacy class is constant. They form a submodule of the space of mappings from \(G\) to \(K\) , denoted by \(\mathcal{Z}_K(G)\) . When \(G\) is finite, the K-algebra \(\mathcal{Z}_K(G)\) is a free K-module of dimension the number of conjugacy classes of \(G\) . The center \(Z(K[G])\) of the algebra \(K[G]\) consists of the elements \(a = \sum a_g g\) of \(K[G]\) such that \(hah^{-1} = a\) for every \(h \in G\) . Now, we have

\[
h a h ^ {- 1} = \sum_ {g \in \mathrm{G}} a _ {h ^ {- 1} g h} g;
\]

consequently, when G is finite, the center \(Z(K[G])\) of the algebra K{[}G{]} consists of the central functions.

Let \((\mathrm{M},\pi)\) be a linear representation of G. Suppose that M is a finite-dimensional free K-module. The trace of \(\pi\) is the trace of the representation of K{[}G{]} associated with \(\pi\) , that is (VIII, p.~382), the linear form \(a\mapsto\mathrm{Tr}(\pi(a))\) on K{[}G{]}. This linear form is determined by the mapping \(g\mapsto\mathrm{Tr}(\pi(g))\) from G to K, which we call the character of the representation \(\pi\) and denote by \(\chi_{\pi}\) . The character of a representation is a central function (II, §4, No.~3, p.~273, Proposition 3).

Let \(M\) and \(M'\) be finite-dimensional free \(K\) -modules. Let \(\pi\) and \(\pi'\) be linear representations of \(G\) in \(M\) and \(M'\) , respectively. Then \(M \oplus M'\) is a finite-dimensional free \(K\) -module, and, by Proposition 1 of III, §9, No.~2, p.~543, we have

\[
\chi_ {\pi \oplus \pi^ {\prime}} = \chi_ {\pi} + \chi_ {\pi^ {\prime}}.
\]

More generally, let

\[
0 \longrightarrow \mathrm{M} ^ {\prime} \longrightarrow \mathrm{M} \longrightarrow \mathrm{M} ^ {\prime \prime} \longrightarrow 0
\]

be an exact sequence of K{[}G{]} modules, and let \(\pi\) , \(\pi'\) , and \(\pi''\) be the linear representations of G associated with M, \(M'\) , and \(M''\) , respectively. Suppose that the K-modules \(M'\) and \(M''\) are free and finite-dimensional. Then so is M (II, §1, No.~11, p.~218, Proposition 21), and we have

\[
\chi_ {\pi} = \chi_ {\pi^ {\prime}} + \chi_ {\pi^ {\prime \prime}}.\tag{1}
\]

PROPOSITION 1. --- Suppose that K is a commutative field. Let \(\pi\) and \(\pi'\) be finite-dimensional semisimple K-linear representations of G.

\begin{enumerate}
\def\labelenumi{\alph{enumi})}
\item
  If the characteristic polynomials of \(\pi(g)\) and \(\pi'(g)\) are the same for every \(g \in G\) , then the representations \(\pi\) and \(\pi'\) are isomorphic.
\item
  Suppose, moreover, that K has characteristic 0. If the characters \(\chi_{\pi}\) and \(\chi_{\pi'}\) are equal, then the representations \(\pi\) and \(\pi'\) are isomorphic.
\end{enumerate}

Assertion a) follows from Corollary 1 of VIII, p.~386; assertion b) follows from part a) of the corollary of VIII, p.~384.

Examples. --- 1) The unit or trivial representation of G is the representation \((\mathrm{K}, \varepsilon)\) , where \(\varepsilon(g) = \operatorname{Id}_{\mathrm{K}}\) for every \(g \in G\) . Its character is the constant function with value 1.

\begin{enumerate}
\def\labelenumi{\arabic{enumi})}
\setcounter{enumi}{1}
\tightlist
\item
  The (left) regular representation of G is the representation \(\gamma\) of G in K{[}G{]} defined by \(\gamma(g)(x)=gx\) for \(g\in G\) and \(x\in K[G]\) . It corresponds to the left regular representation of the algebra K{[}G{]} (VIII, p.~375). Suppose that the group G is finite. The regular representation then has finite degree. For every element g of G different from the identity, the matrix of left multiplication by g in K{[}G{]} with respect to the canonical basis is the matrix of a permutation without fixed point. We therefore have
\end{enumerate}

\[
\chi_ {\boldsymbol {\gamma}} (g) = \left\{ \begin{array}{l l} | \mathrm{G} | & \text { if   g   is   the   identity   element }, \\ 0 & \text { otherwise }. \end{array} \right.\tag{2}
\]

We define the right regular representation of G likewise. The biregular representation of G is the representation \(\rho\) of G × G in K{[}G{]} defined by \(\rho(g, g')(x) = gxg'^{-1}\) for \((g, g') \in \mathrm{G} \times \mathrm{G}\) and \(x \in \mathrm{K}[\mathrm{G}]\) .

\begin{enumerate}
\def\labelenumi{\arabic{enumi})}
\setcounter{enumi}{2}
\item
  Given a linear representation \((\mathrm{M},\pi)\) of G, the contragredient or dual representation \(\pi^{\vee}\) of \(\pi\) is the representation of G in the K-module dual to M defined by the relation \(\pi^{\vee}(g)={}^{t}\pi(g^{-1})\) for every \(g\in G\) (cf.~II, §2, No.~5, p.~234). If M is a finite-dimensional free K-module, then so is its dual, and we have \(\chi_{\pi^{\vee}}(g)=\chi_{\pi}(g^{-1})\) for every \(g\in G\) .
\item
  Let \((\mathrm{M},\pi)\) and \((\mathrm{M}^{\prime},\pi^{\prime})\) be linear representations of G. In Example 1 of VIII, p.~198, we defined a K{[}G{]}-module structure on \(M\otimes_{K}M^{\prime}\) . The corresponding linear representation is called the tensor product of \(\pi\) and \(\pi^{\prime}\) and is denoted by \(\pi\otimes\pi^{\prime}\) . For \(g\in G\) , \(x\in M\) , and \(x^{\prime}\in M^{\prime}\) , we have \(\big((\pi\otimes\pi^{\prime})(g)\big)(x\otimes x^{\prime})=\pi(g)x\otimes\pi^{\prime}(g)x^{\prime}\) .
\end{enumerate}

If M and M' are finite-dimensional free K-modules, then by Proposition 2 of III, §9, No.~2, p.~543, we have

\[
\chi_ {\pi \otimes \pi^ {\prime}} = \chi_ {\pi} \chi_ {\pi^ {\prime}}.\tag{3}
\]

\begin{enumerate}
\def\labelenumi{\arabic{enumi})}
\setcounter{enumi}{4}
\tightlist
\item
  Suppose that G is the product \(G' \times G''\) of two groups. The K-linear mapping from \(K[G'] \otimes_{K} K[G'']\) to K{[}G{]} that sends \(g' \otimes g''\) to \((g', g'')\) for \(g' \in G'\) and \(g'' \in G''\) is an algebra isomorphism. Let \((M', \pi')\) be a linear representation of \(G'\) and \((M'', \pi'')\) a linear representation of \(G''\) . The external tensor product of \(\pi'\) and \(\pi''\) , denoted by \(\pi' \boxtimes \pi''\) , is the representation of \(G' \times G''\) in the vector space \(M' \otimes M''\) defined by \((\pi' \boxtimes \pi'')(g', g'') = \pi'(g') \otimes \pi''(g'')\) for \((g', g'') \in G' \times G''\) . If \(M'\) and \(M''\) are finite-dimensional free K-modules, then \(M' \otimes_{K} M''\) is a finite-dimensional free K-module and the character of the representation \(\pi' \boxtimes \pi''\) is given by the formula
\end{enumerate}

\[
\chi_ {\pi^ {\prime} \boxtimes \pi^ {\prime \prime}} (g ^ {\prime}, g ^ {\prime \prime}) = \chi_ {\pi^ {\prime}} (g ^ {\prime}) \chi_ {\pi^ {\prime \prime}} (g ^ {\prime \prime})
\]

for \(g' \in G'\) and \(g'' \in G''\) (Proposition 2 of III, §9, No.~2, p.~542).

\begin{enumerate}
\def\labelenumi{\arabic{enumi})}
\setcounter{enumi}{5}
\tightlist
\item
  Let \((V,\pi)\) be a linear representation of G such that V is a finite-dimensional free K-module. We define a representation \(\rho\) of \(G\times G\) in \(\operatorname{End}_{K}(V)\) by the formula
\end{enumerate}

\[
\rho (g, g ^ {\prime}) (u) = \pi (g ^ {\prime}) \circ u \circ \pi (g ^ {- 1}).
\]

The K-module isomorphism \(\theta_{\mathrm{V}}: \mathrm{V}^{*} \otimes_{\mathrm{K}} \mathrm{V} \to \mathrm{End}_{\mathrm{K}}(\mathrm{V})\) of II, §4, No.~2, p.~271 is an isomorphism of representations from the external tensor product \(\pi^{\vee} \boxtimes \pi\) to \(\rho\) . The mapping \(g \mapsto \rho(1, g)\) (resp. \(g \mapsto \rho(g, 1)\) ) is a representation of G isomorphic to \(\pi^{\dim_{\mathrm{K}} \mathrm{V}}\) (resp. \((\pi^{\vee})^{\dim_{\mathrm{K}} \mathrm{V}}\) ).

Let L be a commutative K-algebra, and let \((\mathrm{M},\pi)\) be a linear representation of the group G. The group homomorphism \(\pi_{(\mathrm{L})}:\mathrm{G}\to\mathbf{GL}(\mathrm{M}_{(\mathrm{L})})\) defined by \(g\mapsto\mathrm{Id}_{\mathrm{L}}\otimes\pi(g)\) is a linear representation of G in the L-module \(\mathrm{M}_{(\mathrm{L})}\) , called the linear representation of G deduced from the representation \(\pi\) by extension of the ring of scalars K to L.

Suppose that K is a field and that L is a nonzero commutative K algebra. Let \((M,\pi)\) and \((M',\pi')\) be linear representations of G. The representations \(\pi\) and \(\pi'\) are isomorphic if and only if \(\pi_{(L)}\) and \(\pi'_{(L)}\) are (VIII, p.~37, Theorem 3).

Suppose, moreover, that the algebra L is an extension of K. Consider the rings \(\mathrm{R}_{\mathrm{K}}(\mathrm{G})\) and \(\mathrm{R}_{\mathrm{L}}(\mathrm{G})\) defined in Example 1 of VIII, p.~198. Extension of scalars defines a ring homomorphism

\[
u: \mathrm{R} _ {\mathrm{K}} (\mathrm{G}) \longrightarrow \mathrm{R} _ {\mathrm{L}} (\mathrm{G}).
\]

This homomorphism is injective, and an element \(\xi \in \mathrm{R}_{\mathrm{K}}(\mathrm{G})\) is effective if and only if \(u(\xi)\) is (VIII, p.~195, Theorem 1).

